\subsubsection{Pentadic representation of the microfiber and its incidence}
\label{pi-estrella:desarrollo}

The arrangement of the five regions makes it possible to visualize
simultaneously the unity of the ternary quotient and the differentiation
preserved by the catalog. The relevant map is the restriction
\[
 \Psi\big|_{\mathscr R_\pi^{\rm micro}}:
 \mathscr R_\pi^{\rm micro}\longrightarrow\{010211\},
 \qquad \Psi(U)=(-U)\bmod3.
\]
Its five preimages are the vectors \(U_6\) in the preceding table.
Figure \ref{pi-estrella:figura} recovers their radial organization while
preserving the chart order, signatures, and multiplicities of the catalog
constructed in the preceding sections.

% Chart order and data recovered from fig_cinco_regiones_pi.tex.
\begin{figure}[H]
\centering
\begin{tikzpicture}[x=1cm,y=1cm,font=\small,>=Latex]
 \coordinate (a) at (90:2.2);
 \coordinate (b) at (18:2.2);
 \coordinate (c) at (-54:2.2);
 \coordinate (d) at (-126:2.2);
 \coordinate (e) at (162:2.2);
 \draw[black!35,line width=.45pt]
    (a)--(c)--(e)--(b)--(d)--cycle;
 \node[draw,fill=white,circle,minimum size=1.6cm,align=center,
    inner sep=3pt] (z) at (0,0) {$w_\pi$\\$010211$};
 \foreach \v in {a,b,c,d,e}{
    \fill (\v) circle (2pt);
    \draw[->,line width=.6pt] (\v)--(z);
 }
 \node[above=3mm,align=center] at (a)
   {$U_\pi^\perp$\\$555555$\\multiplicity $432$};
 \node[right=3mm,align=left] at (b)
   {$U_{\pi,1}^{++}$\\$339555$\\multiplicity $144$};
 \node[below right=2mm and 1mm,align=left] at (c)
   {$U_{\pi,2}^{++}$\\$393555$\\multiplicity $144$};
 \node[below left=2mm and 1mm,align=right] at (d)
   {$U_{\pi,3}^{++}$\\$933555$\\multiplicity $144$};
 \node[left=3mm,align=right] at (e)
   {$U_\pi^{--}$\\$933717$\\multiplicity $144$};
 \node[align=center,font=\footnotesize] at (0,-3.55)
   {$1_\perp+3_{++}+1_{--}$\qquad $432+4\cdot144=1008$};
\end{tikzpicture}
\caption{The five regions of \(\pi\) in a pentadic arrangement.
Each position preserves its multiplicative signature and its multiplicity;
the radial arrows represent the common projection
\(\Psi(U)=010211\). The star-shaped outline organizes the five chart
positions. Their realization in the five Witt octads is carried out through
the regional alignment and the \(4+1\) decomposition in the text.}
\label{pi-estrella:figura}
\end{figure}


Projection to the center identifies the five ternary images; the record of
the outer positions retains the regional information that makes the
incidence realization possible. The distribution
\(1_\perp+3_{++}+1_{--}\) specifies one transverse region, three positive
regions, and one negative region. Their normalized weights are \(3/7\) and
\(1/7\) for each of the remaining four: the transverse multiplicity is
threefold, even though the five images under \(\Psi\) coincide.

\Needspace{11\baselineskip}
The link with the Witt star is formulated through the already constructed
flag \(B_K\subset H^-_\alpha\), alignment \(\ell\), and lift \(\iota_D\).
The five regional positions are realized in the five octads over
\(\ell(B_K)\):
\[
\begin{aligned}
 U_\pi^\perp&\longmapsto O^\perp_{\ell(B_K)},\\
 U_\pi^{--}&\longmapsto\iota_D(H^-_\alpha),\\
 \{U_{\pi,1}^{++},U_{\pi,2}^{++},U_{\pi,3}^{++}\}
 &\longmapsto
 \{\iota_D(B_K\cup P):P\in\mathcal P_+\},\\
 \mathcal P_+&=\bigl\{\{1,3\},\{2,11\},\{8,12\}\bigr\}.
\end{aligned}
\]
The assignment of the three positive positions to the three pairs is carried
out in an ordered chart. The transverse row and the negative row are
distinguished by their regional predicates and by the flag.

The two figures are thereby articulated: Figure \ref{exc:fig-estrella} shows
the four hexads of the small Witt design over \(B_K\); Figure
\ref{pi-estrella:figura} shows the five regions that, under the alignment,
occupy the four residual octads and the transverse octad of the binary design.
The \(4+1\) law, proved in Proposition \ref{exc:residuo-binario}, specifies
the change of incidence. In this composition, \(K\) determines the marked
tetrad, and the signature of \(\alpha\) distinguishes one of its hexads. In
addition to their common ternary value, the regions preserve the structure
required for this transport to the exceptional realization.
